\begin{center}\begin{tikzpicture}[>=Latex,font=\scriptsize,scale=.83,transform shape]
\tikzset{n/.style={draw=newsblue,rounded corners=2pt,align=center,fill=blue!4,minimum width=21mm,minimum height=8mm}}
\node[n] (a) at (0,0) {中红外吸收\\局部温升};\node[n] (b) at (3,0) {可见光干涉\\折射率变化};\node[n] (c) at (6,0) {多角度数据\\三维重建};
\draw[->,thick] (a)--(b);\draw[->,thick] (b)--(c);
\node[draw=orange!80!black,fill=orange!12,rounded corners,align=center] at (3,-1.1) {19.2个三维体／秒　约52毫秒／体};
\end{tikzpicture}\par{\scriptsize\color{newsgray}图4　MIP-ODT测量流程，\textbf{非原论文仪器结构图}。速度来自Fukushima等，《自然·通讯》，2026年10月30日；编辑部示意。}\end{center}
